\section{Introdução}

\pendente{Contextualizar a fase preparatória das contratações públicas
municipais sob a Lei nº 14.133/2021 e apresentar o problema a partir de fontes
verificáveis. Estrutura pretendida: contexto → problema → justificativa →
organização do trabalho.}

\subsection{Contextualização}

A Lei nº 14.133/2021 organiza o processo de contratação pública em fases, entre
elas a fase preparatória, na qual se inserem o Estudo Técnico Preliminar (ETP),
o Termo de Referência (TR) e a pesquisa de preços que fundamenta o valor
estimado da contratação \cite{brasil2021lei14133}.

\pendente{Sustentar com fonte cada afirmação sobre a duração e o custo dessa
fase em municípios de pequeno porte. Não usar número sem origem — a percepção
de quem trabalha no setor entra como evidência qualitativa levantada em
entrevista, registrada no diário de requisitos, e não como estatística.}

\subsection{Problema de pesquisa}

\pendente{Formular a pergunta em uma frase, no formato "em que medida ...".
Rascunho a validar com a orientadora: em que medida a geração assistida de
minutas e a automação da pesquisa de preços a partir de dados públicos reduzem
o esforço da fase preparatória em uma prefeitura de pequeno porte, sem
comprometer a conformidade legal dos documentos produzidos?}

\subsection{Justificativa}

\pendente{Articular três eixos: relevância prática (gargalo real em municípios
do interior), relevância acadêmica (aplicação de recuperação de informação e
modelagem de dados a um domínio pouco explorado) e viabilidade (acesso a
especialista do domínio e a dados públicos).}

\subsection{Organização do trabalho}

\pendente{Escrever por último, quando as seções estiverem estáveis.}
